\section{A collision experiment on a finite-smoothness class}
\label{sec:setting}
The local identity behind this paper is a physical density formula, not
a boundary-value measurement. For a selected three-impact return,
\[
 f_T(s,t)=\frac{-W_{st}(s,t)}{2\pi A}(T-W(s,t))_+.
\]
Two active densities determine the stationary action $W$ and the free
cell area $A$. Earlier local recognition used certified boundary and
clearance queries to attach the same intrinsic record to different
physical copies. We now construct the needed geometry from additional
short launches. This changes the acquisition cost and the spatial extent
of the local experiment, while keeping the intrinsic law and its absolute
normalization. In particular, a complete boundary is sampled rather than
continued from a short analytic arc.

\begin{definition}[Numerical prior and directed bounds]\label{def:prior}
Fix $0<\beta\le1$. A table consists of
\[
 \mathcal O=\bigcup_{i=1}^r\bigcup_{\lambda\in\Lambda}(C_i+\lambda),
 \qquad A=\covol\Lambda-\sum_{i=1}^r\area(C_i).
\]
The prior supplies numerical upper bounds
\[
 r_0,D_0,A_1,V_1,b,L_0,M_6,\kappa_+,
\]
and positive lower bounds
\[
 A_0,V_0,a_0,d_0,\kappa_-,\eta,\sigma.
\]
In particular,
\[
 r\le r_0,\quad \operatorname{diam}C_i\le D_0,\quad
 A_0\le A\le A_1,\quad V_0\le\covol\Lambda\le V_1,\quad
 \area(C_i)\ge a_0.
\]
Different physical bodies have distance at least $d_0$. Their curvature
lies in $[\kappa_-,\kappa_+]$. Some presentation $\Lambda=L\Z^2$ has
$\|L\|,\|L^{-1}\|\le L_0$ and a fundamental parallelogram of diameter
at most $b$. No such presentation is supplied. Centered support functions
$p_i$ satisfy $\|p_i\|_{C^{6,\beta}}\le M_6$ and
\begin{align}
 \|p_i(\cdot+\phi)-p_i\|_2&\ge\eta\dist(\phi,2\pi\Z),\label{eq:rot}\\
 \|p_i(-\cdot+\phi)-p_i\|_2&\ge\eta,\label{eq:reflection}\\
 \inf_{O\in O(2)}\|p_i-p_{OC_j}\|_2&\ge\sigma\quad(i\ne j).\label{eq:types}
\end{align}
Choose $2b<R_-<R_+$. Every clear shortest bridge of length below $R_-$
has fixed positive clearance, contact-chart, adjacent-flight, twist,
non-grazing and time-collar margins. Additional admitted bridges have
length at most $R_+$ and the same specified smaller margins. Numerical
lower bounds for these margins are denoted $c_0,r_c,\ell_0,\omega_0,v_*,d_*$,
respectively; numerical bounds for the required chart derivatives are
included. There are fixed nested protected endpoint rectangles on which
$W-2g\le d_*/12$. Active validation times have fixed positive separation.
All bounds are directed: an upper bound may overestimate the true value;
a lower bound may underestimate it, but must be strictly positive.
\end{definition}

One may work in a closed subclass satisfying these bounds with slack
where a protected neighborhood is used. Compactness needs no analytic
strip. Keep bounded presentations $L$ and representative centers $Lt_i$,
$t_i\in[0,1]^2$, instead of choosing a discontinuous reduced basis.
The $C^{6,\beta}$ bound is compact in $C^6$; finitely many body counts
and bounded integer bridge labels suffice on bounded neighborhoods.
Changes at representative faces merely retain another presentation of
the same table. Thus a closed class is compact in the lower norms used
below. The finite-smoothness hypothesis implies uniform $C^{2,\beta}$
endpoint densities: arclength boundary charts are $C^{5,\beta}$,
stationary reduction preserves that regularity for $W$ by the envelope
identity, and the density differentiates $W$ only twice.

The margins in \eqref{eq:rot}--\eqref{eq:types} are still substantive.
They imply that $\Lambda$ is the full translation group: a period must
carry a body to a congruent type, hence to its own declared translate,
and a compact body has no nonzero translational symmetry. This paper
does not infer the numerical priors from the observations. Unlike the
preceding calibration, it requires neither complex-strip bounds nor
local holomorphic graph radii. Smooth nonanalytic perturbations are allowed.

\subsection{The launch and localization contract}\label{sec:sensor}
The apparatus can choose a position uniformly in a specified laboratory
square and choose an independent uniform direction at speed one. It
reports a solid initial position as a failure. At a free position it
records the ordered collisions up to a prescribed bounded time, with
laboratory position enclosures of radius at most a requested tolerance
$\varepsilon_x$. Short auxiliary launches may use translated squares
around the observed contacts and candidate segment. These squares have
a prior-dependent bounded size, large enough to cover an encountered
body of diameter $D_0$, not merely an infinitesimal contact patch.
All launches, including solid positions and trajectories with no collision,
are charged. Position-localization work, travel and setup are separate
from the number of launches. The collision-time and launch controls are
calibrated; no boundary height, derivative, normal, arclength, body label,
period, or distance-to-solid query is supplied.

This is a resettable spatial experiment. Calling it a passive unmarked
trajectory would be incorrect. Its geometric outputs below are confidence
enclosures, with explicitly spent error probability, not deterministic
oracles. The final two windows are two growing two-dimensional endpoint
histograms. They are not two scalar observations.

\begin{theorem}[Collision-generated local geometry]\label{thm:probe-main}
Under the diameter, separation, curvature and $C^3$ consequences of
Definition~\ref{def:prior}, there are computable $C,c,\nu_0>0$ such that
at an encountered bounded flight, $0<\nu<\nu_0$ and $0<\zeta<1/4$,
at most
\begin{equation}\label{eq:probe-main}
 C\nu^{-3/2}\log\frac{C}{\nu\zeta}
\end{equation}
independent auxiliary launch attempts, localized to
$\varepsilon_x\le c\nu^3$, give, with probability at least $1-\zeta$:
complete $C^2$ support-function approximations of the encountered bodies,
local intrinsic coordinates to error $C\nu$, and distance-to-solid
enclosures on a prescribed bounded flight neighborhood. A protected clear
normal bridge is recognized by finite tests with slack. No unknown body
catalogue or analytic continuation is an input to this procedure.
The probability guarantee is uniform conditional on the preceding
history, including the location of the flight.
\end{theorem}

\begin{theorem}[Geometric-size fingerprints]\label{thm:fingerprint-main}
Let
\begin{gather}
 K_0=R_++D_0+1,\qquad C_0=3+9(R_++2D_0)/\eta,\nonumber\\
 \Delta_0=\min\{1,\sigma/36,\eta/36,d_0/(4C_0)\},\quad
 \chi=\Delta_0/4,\quad m=2\lceil32K_0/\Delta_0\rceil.\label{eq:fingerprint-size}
\end{gather}
The gap and the support values of both complete bodies, in their common
normal-contact frame at $m$ equally spaced angles, separate all oriented
bounded bridge classes coexisting in one table by more than $\chi$.
Simultaneous transverse reversal permutes these coordinates.
Theorem~\ref{thm:probe-main} acquires this descriptor to any prescribed
smaller error. Its dimension $1+2m$ is polynomial in the numerical bounds
and reciprocal margins appearing in \eqref{eq:fingerprint-size}; it has
no analytic-continuation exponent. This is not a polynomial bit-complexity
claim and does not identify arbitrary bodies from different tables.
\end{theorem}

\begin{theorem}[Smooth whole-table certification]\label{thm:main}
On Definition~\ref{def:prior}'s class, using only the launch and localization
contract above, a sequential procedure has the following properties.
With probability at least $1-\delta$, $0<\delta<1/4$, every certificate
issued at any epoch is valid for the entire table and its full period
group, up to one common Euclidean isometry, permutation, representative
changes and an integer basis change. Neither a witness graph nor the
number and identities of the types are supplied.

For a chosen $a\ge6$ and target accuracy $\xi>0$, let $T_\xi$ be the
first passing epoch. There are computable class constants $p_*>0,w_0<\infty$
and a deterministic $k_\xi$ such that
\begin{equation}\label{eq:tail-main}
 \Prob(T_\xi>k)\le \frac{\delta}{4(k+1)^a}+w_0e^{-p_*k}
 \quad(k\ge k_\xi).
\end{equation}
In particular $T_\xi<\infty$ almost surely. Writing
\[
 r_k=\frac{\log(C(k+1)^{a+6}/\delta)}{(k+1)^2},
\]
the certified error is at most $C r_k^{\beta/(2\beta+6)}$ for all
sufficiently large passing epochs. No continuation power is applied.
Through epoch $m$ the total number of attempts, including auxiliary
geometric launches, is at most
\begin{equation}\label{eq:cost-main}
 C(m+1)^{11/2}\log\frac{C(m+1)}{\delta}
\end{equation}
with the separately stated rejection-tail probability. Its expectation
up to $T_\xi$ is finite, since $a>11/2$. The largest square half-width is
$O((m+1)^{(\beta+4)/(\beta+3)})$. Localization precision, geometric
postprocessing, travel and setup are not hidden in the launch count.
\end{theorem}

The proof first constructs nearby physical geometry from collision clouds,
then uses the original intrinsic density normalization to exclude missing
types or an unsaturated period subgroup. A finite-smoothness boundary
cannot be inferred from one of its short arcs; no such inference is made.
The full body is sampled in a bounded region. The earlier analytic
short-patch theorem remains a different valid route in the retained volume.
